\begin{thm}[OPAC-034] \label{thm-opac-034}
    Every Banff quiver is mutation-equivalent to a quiver in $\BB'$, \ie $\BB = \BB'$.
\end{thm}

\begin{proof}
    We argue by induction on the number of vertices in a quiver.
    The single vertex quiver is acyclic, and thus it is in both $\BB$ and $\BB'$, completing the base case.
    Suppose then that every Banff quiver on $n$ vertices is contained in $\BB'$.
    Let $Q$ be a Banff quiver on $n+1$ vertices.
    As both $\BB$ and $\BB'$ are closed under mutation, we can assume that $Q$ has a covering pair $(i,j)$ such that $Q \setminus \{i\}$ and $Q \setminus \{j\}$ are Banff.
    If either $i$ is a source or $j$ is a sink, then $Q \in \BB'$ by the inductive hypothesis.
    We can then assume that $i$ is not a source and $j$ is not a sink and split into the two cases given by Proposition~\ref{prop-acyclic-covering-pair}.

    Then Lemma~\ref{lem-covering-pair-mutation} gives us a source or sink mutation-sequence $\bw$ such that either $i$ is a source or $j$ is a sink in $P = \mu_\bw(Q)$.
    Additionally, we get that $P \setminus \{i\}$ and $P \setminus \{j\}$ are mutation-equivalent to $Q \setminus \{i\}$ and $Q \setminus \{j\}$.
    Both $P \setminus \{i\}$ and $P \setminus \{j\}$ are then Banff and in $\BB'$ by the inductive hypothesis.
    This implies that $P \in \BB'$, forcing $Q \in \BB'$.
    As $Q$ was an arbitrary member of $\BB$, we have $\BB = \BB'$ by induction.
\end{proof}
